\chapter{Kết luận và hướng phát triển}
\label{chap:ket-luan}

\section{Kết luận}

Khoá luận đã xây dựng AI-CORP, một nền tảng đa tác tử dùng LLM, mô phỏng một công ty phần mềm có tám vai trò: CEO, BA, Architect, PM, DEV, QC, QA và Marketing. Phân tích và thiết kế hướng đối tượng được áp ở ba cấp, gồm kiến trúc của nền tảng, thiết kế của từng tác tử, và sản phẩm mà các tác tử làm ra cho người dùng.

Ở cấp kiến trúc, tầng miền không phụ thuộc framework. Mỗi tác tử dựng theo Template Method cùng hai Strategy. Tác tử, công việc và lượt chạy đều có máy trạng thái riêng, các mô-đun nói chuyện với nhau qua sự kiện miền, và mọi phụ thuộc vào hạ tầng đi qua Repository hoặc cổng. Bộ máy quy trình hỗ trợ quy trình con, phiên bản bất biến, ảnh chụp và khôi phục, nhánh song song, phê duyệt của con người và ba kiểu cộng tác giữa tác tử. Ở cấp sản phẩm, dự án nào cũng đi qua use case, mô hình miền, sơ đồ UML, mã bốn tầng, kiểm thử chấp nhận và ma trận truy vết, với cổng OOAD, phân tích tĩnh và một vòng tự sửa có giới hạn chi phí canh ở giữa. Ở cấp vận hành, nền tảng có phân quyền theo vai trò, nhật ký kiểm toán, khoá API, webhook ký HMAC, xuất báo cáo, quản lý nhà cung cấp và mô hình LLM ngay trong ứng dụng, điều khiển lượt chạy có ràng buộc ngân sách, phân tích hiệu năng tác tử, dự báo thời gian và chi phí, cùng cấu hình triển khai Docker có kiểm tra cấu hình production và điểm kiểm tra tình trạng. Giao diện web được làm lại trên Ant Design 6, có chế độ sáng, tối và chế độ gọn, mà không thay đổi logic hay API.

Những kết quả chính đo được là:

\begin{itemize}
\item 707 kiểm thử API và 44 kiểm thử web đều đạt, tầng miền có độ bao phủ 97,1\% số dòng và 88,1\% số nhánh.
\item Ba lần chạy đầu cuối trên hạ tầng thật đạt 24/24, 21/21 và 44/44 kiểm tra. Cấu hình production dựng được lược đồ từ migration và không lệch so với mô hình.
\item Mọi nút quy trình xong dưới 2 giây (p95 là 1,64~s, chưa trừ độ trễ LLM giả lập). Một tiến trình API phục vụ khoảng 605 yêu cầu đọc mỗi giây với 20 người dùng và bão hoà ở khoảng 407 yêu cầu mỗi giây với 50 người dùng, không có lỗi.
\item Dự báo thời gian tác tử có MAPE leave-one-out từ 3\% đến 30\% tuỳ vai trò.
\item Kiểm thử trên hạ tầng thật, rà soát giao diện và logic, cùng việc cài đặt mới đã tìm và sửa 50 lỗi, trong đó có lỗ hổng cho phép người xem phê duyệt qua WebSocket và lỗi cho phép một phê duyệt được trả lời hai lần.
\end{itemize}

Về ba câu hỏi nghiên cứu, có thể trả lời ngắn gọn như sau. Với RQ1, OOAD làm tính mô-đun, khả năng mở rộng và khả năng kiểm thử tốt lên rõ, miễn là có kiểm thử đầu cuối cho thấy thiết kế thật sự được dùng. Với RQ2, các mẫu hiệu quả nhất là Template Method cùng Strategy cho tác tử, Strategy tự đăng ký cùng Composite cho quy trình, và State gắn đúng vào thứ có vòng đời. Nhìn chung, mẫu thiết kế có giá trị nhất khi nó tạo ra một điểm kiểm soát duy nhất có tính nguyên tử. Với RQ3, học từ lịch sử để dự báo thì đo được và cho sai số chấp nhận được. Còn học để làm sản phẩm tốt hơn thì đã có cơ chế và cơ chế chạy đúng, nhưng muốn đo khách quan mức cải thiện thì phải dùng LLM thật trong một thí nghiệm có nhóm đối chứng.

\section{Hướng phát triển}

\begin{enumerate}
\item \textbf{Đánh giá với LLM thật.} Trang quản lý nhà cung cấp đã cho phép chọn mô hình thật cho từng vai trò mà không sửa mã, nên vướng mắc chính chỉ còn là tiền và thời gian chạy. Việc cần làm là chạy lại tập dự án đánh giá với một hoặc vài mô hình thật, đo tỉ lệ đạt cổng OOAD ngay lần đầu, số lần QA từ chối, số kiểm thử chấp nhận đạt, thời gian và chi phí thực tế, rồi đánh giá lại mô hình dự báo và bộ dự đoán chất lượng mã trên dữ liệu mới.
\item \textbf{Thí nghiệm có đối chứng cho cơ chế học}, theo thiết kế ở chương~\ref{chap:thao-luan}, với hai nhóm bật học và tắt học, lặp nhiều lần, có kiểm định thống kê.
\item \textbf{Chuyên gia chấm thiết kế.} Mời chuyên gia chấm sản phẩm OOAD của các dự án sinh ra theo một thang đo, để xem cổng chất lượng tự động tương quan với đánh giá của con người tới đâu.
\item \textbf{Đưa sandbox container vào kiểm thử thường xuyên}, tức là chạy kiểm thử sandbox trên máy CI có Docker và kiểm chứng đầu cuối luồng thực thi mã trong container.
\item \textbf{Mở rộng ngang}, bằng cách thêm adapter Redis cho Socket.IO, chạy nhiều bản sao API sau nginx và thử tải với cấu hình đó.
\item \textbf{Hỗ trợ ngôn ngữ đích khác}, bằng cách tách bộ phân tích tĩnh và trình chạy kiểm thử thành adapter theo ngôn ngữ, đặt sau cổng \code{ICodeAnalyzer} sẵn có, bắt đầu với Python hoặc Java.
\item \textbf{Tăng vai trò của con người}, ví dụ cho phép sửa trực tiếp một sản phẩm, ví dụ mô hình miền, ngay tại bước phê duyệt rồi để quy trình chạy tiếp từ bản đã sửa, tận dụng cơ chế ảnh chụp và khôi phục đã có.
\item \textbf{Giao diện đa ngôn ngữ và trợ năng}: thêm bản tiếng Việt cho nhãn giao diện, kiểm tra với trình đọc màn hình và trên màn hình nhỏ.
\item \textbf{Kiểm thử xâm nhập độc lập} cho API, WebSocket và webhook trước khi dùng với dữ liệu thật.
\item \textbf{Hoàn thiện điều khiển chi phí và phiên làm việc}, gồm huỷ lời gọi LLM quá hạn, ước lượng chi phí trước mỗi lời gọi để không nút nào làm vượt ngân sách, cho phép giao lại phê duyệt, thêm nhánh riêng khi bị từ chối và hỗ trợ token làm mới.
\end{enumerate}
